%% =====================================================================
%%  Prague Cafe Project -- report, topic: Cafe Size
%%  Tables and figures are produced by the R scripts in ../analysis
%%  and loaded here; captions are written here, not in R.
%% =====================================================================
\documentclass[11pt,a4paper]{article}
\usepackage{cafestyle}
\usepackage{adjustbox}
\usepackage{array}
\usepackage{placeins}   % \FloatBarrier keeps each table and figure in its own section
% P{width}: a table column whose text wraps and is left-aligned (used by the R tables)
\newcolumntype{P}[1]{>{\raggedright\arraybackslash}p{#1}}

% Include an R-exported file if it exists, otherwise show a visible placeholder
\newcommand{\inputtable}[1]{\IfFileExists{#1}{\input{#1}}{\par\noindent\fbox{\parbox{0.9\linewidth}{\centering Missing table file: \texttt{\detokenize{#1}}}}\par}}
% a table at small size, scaled down only if it is still wider than the page
\newcommand{\fittable}[1]{\small\begin{adjustbox}{max width=\linewidth}\inputtable{#1}\end{adjustbox}}
\newcommand{\inputfigure}[2][0.85]{\IfFileExists{#2}{\includegraphics[width=#1\linewidth]{#2}}{\fbox{\parbox{0.8\linewidth}{\centering Missing figure file: \texttt{\detokenize{#2}}}}}}

\title{\textbf{Small caf\'es are rated higher, big caf\'es barely charge more: caf\'e size in Prague}}
\author{Sawyer Gatdula\\[0.3ex]\small Data Analytics, topic: Caf\'e Size\\[0.3ex]\small Partner: Anthony George}
\date{\today}

\begin{document}
\maketitle

\begin{abstract}
\noindent This report looks at how big Prague caf\'es are and whether their size is linked to where they are, how they are rated and what they charge. The data are 189 caf\'es visited by the students of our course. A caf\'e has 28 seats inside on average (median 22), with no evidence that this differs from 30, and about 17 more seats inside than outside. Prague~5 has fewer outdoor seats than Prague~1 and Prague~7. Small caf\'es (15 or fewer seats inside) are top-rated 28.6 percentage points more often than other caf\'es. Size and price are only weakly linked: the Price Level is 0.66~CZK higher per 10 seats, but this link is no longer significant once the distance from the Orloj is taken into account, and it only shows up in caf\'es without outdoor seating.
\end{abstract}

\FloatBarrier
\section{Introduction}
How big a caf\'e is changes what it feels like to be there. A small caf\'e can feel personal and quiet, while a big one has room for groups and tourists. Size might also be linked to location, ratings and prices, for example if big caf\'es tend to be in expensive, touristy areas. This report is the Caf\'e Size chapter of our class study of the Prague caf\'e market. It asks how many seats caf\'es have inside and outside, whether outdoor seating differs between districts, whether small caf\'es get better ratings, and whether bigger caf\'es charge more, also once their distance from the city centre and their outdoor seating are taken into account.

\FloatBarrier
\section{Data and methods}

\FloatBarrier
\subsection{Data}
The data were collected in person by the ten students of the course in autumn 2026. Each of us visited 20 caf\'es and recorded the seats, prices, Google rating and location in a shared sheet. The raw file has 299 rows. The cleaning script (\texttt{01\_cleaning.R}) made these choices:
\begin{itemize}
  \item 99 empty rows at the bottom of the sheet were removed, which left 200 caf\'es.
  \item 11 caf\'es had been entered twice by different students (same address and almost the same number of Google ratings). I kept the first entry of each, which left 189 unique caf\'es. Caf\'es with the same name but a different address, like the three branches of EMA Espresso Bar, were kept because they are different places.
  \item District labels were cleaned up (typos like ``prauge 1'', capitals and extra spaces). Five caf\'es had a district that did not match their own address, so I corrected them: four Mal\'a Strana caf\'es from Prague~2 to Prague~1, and one Ko\v{s}\'i\v{r}e caf\'e from Prague~7 to Prague~5.
  \item 18 caf\'es listed several prices in one cell (for example ``85, 95, 105''), and I kept the first listed price. Empty cells and ``na'' were set to missing, never to 0.
\end{itemize}
The variables used are the seats inside and outside, outdoor seating (at least one outdoor seat), small caf\'e (15 or fewer seats inside), top-rated (Google rating of 4.7 or higher), the Price Level (the average of the espresso, cappuccino and flat white prices, only for caf\'es that have all three), an above-median Price Level (above 79.67~CZK) and the distance from the Orloj. Missing values: seats inside 1, seats outside 3, rating 3, Price Level 2. No variable was transformed. All descriptive statistics are in Appendix~\ref{app:desc}. A typical caf\'e has a median of 22 seats inside and 6 outside, and the mean Price Level is 80.70~CZK.

\FloatBarrier
\subsection{Statistical methods}
All tests use $\alpha=.05$. Before running any test, the descriptive script (\texttt{02\_descriptive.R}) checked every variable in every compared group with a histogram, a Q--Q plot and its skewness (Appendix~\ref{app:dist}). I counted a variable as normal if its skewness was between $-0.5$ and $0.5$. If it was not, I checked whether the central limit theorem applies: the group needed at least 30 caf\'es, and the means of 5{,}000 bootstrap resamples (seed 123) needed a skewness between $-0.5$ and $0.5$. No variable was normal on its own, but the central limit theorem applied to all of them except the outdoor seats of Prague~2, 3 and 5, which have fewer than 30 caf\'es (Table~\ref{tab:dtable}). So A1--A3 use $t$-tests (Welch's test for A2), A4 uses the Kruskal--Wallis test with Dunn's post hoc test and Holm's correction, and A6 uses Pearson's $r$. A5 uses the chi-square test; Fisher's exact test would have been used if an expected count were below 5. A7--A9 are linear regressions. I checked their residuals after fitting and planned robust (HC3) standard errors for any model where the Breusch--Pagan test showed unequal variance (none did). A10 is a logistic regression, reported as an average marginal effect. Confidence intervals for means come from the $t$ distribution, and those for medians from a percentile bootstrap (5{,}000 resamples, seed 123). The direction of every one-tailed test was fixed before looking at the data. The analysis was done in R 4.3.3 with dplyr 1.1.4, ggplot2 3.4.4, xtable 1.8-4, rstatix 0.7.2, lmtest 0.9-40, sandwich 3.1-0, emmeans 1.10.0 and marginaleffects 0.18.0.

\FloatBarrier
\section{Results}

Table~\ref{tab:summary} summarises all ten analyses. Below, they follow the order of the story: how big caf\'es are (A1--A3), where outdoor seating is (A4), how size relates to ratings (A5) and how it relates to prices (A6--A10).

\begin{table}[ht]
\centering
\caption{Main results of the ten analyses. Estimates with 95\% confidence intervals; $p$-values of A2, A3 and A5 are one-tailed and the direction was fixed in advance. pp = percentage points.}
\label{tab:summary}
\fittable{tables/summary.tex}
\end{table}

\FloatBarrier
\subsection{A1: How many seats inside?}
\Ho: the mean number of seats inside is 30. \Hone: it is different from 30 (two-tailed). The 188 caf\'es with a seat count (1 missing) have \meanci{27.99}{23.78}{32.20}{seats}, a median of 22 and an interquartile range of 17 (Table~\ref{tab:a1}, Figure~\ref{fig:a1}). The difference from 30 is tiny, Cohen's $d=-0.07$, and the one-sample $t$-test finds no evidence that the mean differs from 30, \tstat{187}{$-$0.94}, \pval{.347}. In this sample a caf\'e has about 28 seats inside on average. The median of 22 is lower because a few very big caf\'es, like VNITROBLOCK with 260 seats, pull the mean up.

\begin{table}[ht]
\centering
\caption{A1: seats inside compared with 30. $n$ = caf\'es used, Lost = caf\'es with a missing value.}
\label{tab:a1}
\fittable{tables/A1_table.tex}
\end{table}

\begin{figure}[ht]
\centering
\inputfigure[0.7]{figures/A1_figure.pdf}
\caption{A1: seats inside of each caf\'e (grey dots, $n=188$), the mean (blue point) with its 95\% confidence interval (blue bar) and the tested value of 30 seats (dashed red line).}
\label{fig:a1}
\end{figure}

\FloatBarrier
\subsection{A2: Do caf\'es with outdoor seating have more seats inside?}
\Ho: caf\'es with outdoor seating have the same number of seats inside as caf\'es without, or fewer. \Hone: they have more (one-tailed, fixed in advance because the question asks about ``more''). The difference is caf\'es with outdoor seating minus caf\'es without. The 124 caf\'es with outdoor seating have 30.69 seats inside on average, and the 62 without have 23.06 (3 caf\'es lost to missing values). The difference is 7.62~seats, \CI{$-$0.47}{15.72}, with Cohen's $d=0.27$ (Table~\ref{tab:a2}, Figure~\ref{fig:a2}). Levene's test found no difference in variances, $F(1,184)=0.17$, \pval{.680}, and I used Welch's test, which does not need equal variances anyway. Welch's $t$-test supports \Hone, \tstat{159.48}{1.86}, one-tailed \pval{.032}. So caf\'es with outdoor seating have about 7.6 more seats inside, a small effect ($d=0.27$). But the two-sided interval still includes 0, so how big the difference really is remains uncertain.

\begin{table}[ht]
\centering
\caption{A2: seats inside of caf\'es with and without outdoor seating. Difference = with minus without, with its two-sided 95\% confidence interval; $d$ = Cohen's $d$; the $p$-value is one-tailed. $n$ = caf\'es in each group, Lost = caf\'es with a missing value.}
\label{tab:a2}
\fittable{tables/A2_table.tex}
\end{table}

\begin{figure}[ht]
\centering
\inputfigure[0.7]{figures/A2_figure.pdf}
\caption{A2: seats inside of caf\'es with outdoor seating ($n=124$) and without ($n=62$). Grey dots are caf\'es; blue points and bars are the group means with 95\% confidence intervals.}
\label{fig:a2}
\end{figure}

\FloatBarrier
\subsection{A3: How many more seats inside than outside?}
\Ho: the mean difference, inside minus outside, is 0 or less. \Hone: it is greater than 0 (one-tailed). The test is paired because both counts come from the same caf\'e. The 186 caf\'es with both counts (3 lost) have 28.15 seats inside and 10.78 outside on average. The mean difference is 17.36~seats, \CI{13.39}{21.33}, with $d_z=0.63$ (Table~\ref{tab:a3}, Figure~\ref{fig:a3}). The paired $t$-test supports \Hone, \tstat{185}{8.63}, one-tailed \plt{.001}. In this sample a caf\'e has about 17 more seats inside than outside. That is more than half of its average indoor seating, so the difference clearly matters in practice.

\begin{table}[ht]
\centering
\caption{A3: seats inside and outside of the same caf\'es. Difference = mean of inside minus outside, with its two-sided 95\% confidence interval; $d_z$ = mean difference divided by the SD of the differences; the $p$-value is one-tailed.}
\label{tab:a3}
\fittable{tables/A3_table.tex}
\end{table}

\begin{figure}[ht]
\centering
\inputfigure[0.7]{figures/A3_figure.pdf}
\caption{A3: the difference seats inside minus seats outside of each caf\'e (grey dots, $n=186$), the mean difference (blue point) with its 95\% confidence interval (blue bar), and no difference (dashed red line at 0).}
\label{fig:a3}
\end{figure}

\FloatBarrier
\subsection{A4: Do outdoor seats differ between districts?}
\Ho: the number of outdoor seats is the same in Prague~1, 2, 3, 5, 7 and 8. \Hone: at least one district is different. Prague~2, 3 and 5 have fewer than 30 caf\'es and their counts are skewed, so I used the rank-based Kruskal--Wallis test. For the 176 caf\'es in these districts with an outdoor count (3 lost), $H(5)=12.51$, \pval{.028}, so the districts differ (Table~\ref{tab:a4}, Figure~\ref{fig:a4}). Dunn's post hoc test with Holm's correction (Appendix~\ref{app:checks}, Table~\ref{tab:a4posthoc}) shows that Prague~5, with a median of 0 outdoor seats, has fewer outdoor seats than Prague~1 (difference in medians $-$8 seats, \CI{$-$12.00}{$-$3.00}, Holm \pval{.014}) and Prague~7 (Prague~7 has 9 more, \CI{0.00}{17.50}, Holm \pval{.044}). No other pair differs after the correction. A median of 0 means that at least half of the caf\'es we visited in Prague~5 have no outdoor seating at all.

\begin{table}[ht]
\centering
\caption{A4: outdoor seats by district, medians with bootstrapped 95\% confidence intervals (5{,}000 resamples, seed 123), and the Kruskal--Wallis test.}
\label{tab:a4}
\fittable{tables/A4_table.tex}
\end{table}

\begin{figure}[ht]
\centering
\inputfigure[0.8]{figures/A4_figure.pdf}
\caption{A4: outdoor seats of each caf\'e by district (grey dots, $n=176$). Blue points and bars are the district medians with bootstrapped 95\% confidence intervals.}
\label{fig:a4}
\end{figure}

\FloatBarrier
\subsection{A5: Are small caf\'es more likely to be top-rated?}
\Ho: small caf\'es (15 or fewer seats inside) are top-rated (Google rating 4.7 or higher) no more often than other caf\'es. \Hone: they are top-rated more often (one-tailed). The difference is small caf\'es minus other caf\'es. Of the 186 caf\'es with both values (3 lost), 30 of 52 small caf\'es (57.7\%) and 39 of 134 other caf\'es (29.1\%) are top-rated. That is a difference of 28.6 percentage points, \CI{13.1}{44.1} (Table~\ref{tab:a5}, Figure~\ref{fig:a5}). The smallest expected count is 19.29, above 5, so the chi-square test can be used: \chisqstat{1}{186}{13.12}, one-tailed \plt{.001}, Cram\'er's $V=.27$. The standardised residuals ($\pm$3.62) show that small caf\'es have more top ratings than we would expect if size and rating were unrelated. In this sample a small caf\'e is about twice as likely to be top-rated, which is a big difference in practice.

\begin{table}[ht]
\centering
\caption{A5: top-rated caf\'es (Google rating 4.7 or higher) among small caf\'es (15 or fewer seats inside) and other caf\'es. Top: counts, row percentages and standardised residuals. Bottom: the difference small minus other in percentage points (pp) with its two-sided 95\% confidence interval, Cram\'er's $V$, the smallest expected count and the chi-square test with its one-tailed $p$-value.}
\label{tab:a5}
\fittable{tables/A5_table.tex}

\medskip
\fittable{tables/A5_test.tex}
\end{table}

\begin{figure}[ht]
\centering
\inputfigure[0.7]{figures/A5_figure.pdf}
\caption{A5: share of top-rated caf\'es among small caf\'es ($n=52$) and other caf\'es ($n=134$), with 95\% confidence intervals and the counts.}
\label{fig:a5}
\end{figure}

\FloatBarrier
\subsection{A6: Is the Price Level correlated with the seats inside?}
\Ho: the Price Level and the seats inside are not correlated. \Hone: they are correlated (two-tailed). Both variables passed the central limit theorem check, so I used Pearson's $r$. For the 186 caf\'es with both values (3 lost), $r=.15$, \CI{.01}{.29}, and $r^2=.02$ (Table~\ref{tab:a6}, Figure~\ref{fig:a6}); \tstat{184}{2.11}, \pval{.037}. Bigger caf\'es tend to be a little more expensive, but the link is weak: the seats share only about 2\% of their variation with the Price Level. A correlation also does not show that size causes higher prices.

\begin{table}[ht]
\centering
\caption{A6: Pearson correlation between the seats inside and the Price Level.}
\label{tab:a6}
\fittable{tables/A6_table.tex}
\end{table}

\begin{figure}[ht]
\centering
\inputfigure[0.7]{figures/A6_figure.pdf}
\caption{A6: Price Level (CZK) against seats inside for 186 caf\'es, with Pearson's $r$ and its 95\% confidence interval.}
\label{fig:a6}
\end{figure}

\FloatBarrier
\subsection{A7: By how many CZK does the Price Level change per 10 seats?}
\Ho: the slope of the Price Level on the seats inside is 0. \Hone: it is not 0 (two-tailed). I divided the seats by 10, so the slope is in CZK per 10 seats. For 186 caf\'es (3 lost), the Price Level goes up by 0.66~CZK per 10 seats inside, \CI{0.04}{1.27}, \tstat{184}{2.11}, \pval{.037}, $R^2=.02$ (Table~\ref{tab:a7}, Figure~\ref{fig:a7}). The Breusch--Pagan test found no unequal variance (\pval{.186}), so classic standard errors are reported. The diagnostics (Appendix~\ref{app:diag}) flag 11 influential caf\'es (Cook's distance above $4/n$), mainly VNITROBLOCK (260 seats), Caf\'e Louvre (200 seats), Starbucks Reserve and Caf\'e Slavia. Without them the slope goes up to 1.92~CZK per 10 seats, so the positive link does not depend on them. In practice the effect is very small: 10 extra seats go with less than 1~CZK more on a Price Level of about 80~CZK.

\begin{table}[ht]
\centering
\caption{A7: linear regression of the Price Level (CZK) on the seats inside (per 10 seats), classic standard errors.}
\label{tab:a7}
\fittable{tables/A7_table.tex}
\end{table}

\begin{figure}[ht]
\centering
\inputfigure[0.7]{figures/A7_figure.pdf}
\caption{A7: Price Level (CZK) against seats inside for 186 caf\'es, with the fitted regression line and its 95\% confidence band.}
\label{fig:a7}
\end{figure}

\FloatBarrier
\subsection{A8: Does the link hold once distance from the Orloj is included?}
\Ho: with the distance from the Orloj held fixed, the slope of the seats is 0. \Hone: it is not 0 (two-tailed). I refitted the A7 model on the same 186 caf\'es with the distance added (Table~\ref{tab:a8}, Figure~\ref{fig:a8}). The slope of the seats dropped from 0.66 to 0.58~CZK per 10 seats, \CI{$-$0.03}{1.20}, \tstat{183}{1.87}, \pval{.063}, so there is no longer evidence that it differs from zero. The distance itself does matter: each km farther from the Orloj goes with a Price Level 1.48~CZK lower, \CI{$-$2.85}{$-$0.12}, \pval{.033}, and adding it improves the model, \Fstat{1}{183}{4.61}, \pval{.033}, $R^2=.05$. The two predictors barely overlap (VIF = 1.01), and the Breusch--Pagan test found no unequal variance (\pval{.152}). Most of the seats slope is still there, but once location is included, the data no longer give evidence that size by itself goes with price.

\begin{table}[ht]
\centering
\caption{A8: the model of A7 and the model with the distance from the Orloj, on the same 186 caf\'es; the $F$-test compares the two models.}
\label{tab:a8}
\fittable{tables/A8_table.tex}
\end{table}

\begin{figure}[ht]
\centering
\inputfigure[0.7]{figures/A8_figure.pdf}
\caption{A8: Price Level (CZK) against seats inside (grey dots, $n=186$), with the fitted line of A7 (seats only) and of A8 (seats and distance, drawn at the median distance), each with its 95\% confidence band.}
\label{fig:a8}
\end{figure}

\FloatBarrier
\subsection{A9: Does the link differ with and without outdoor seating?}
\Ho: the slope of the Price Level on the seats is the same for caf\'es with and without outdoor seating. \Hone: the slopes differ (two-tailed). For 184 caf\'es (5 lost), the interaction is $-$2.67~CZK, \CI{$-$4.17}{$-$1.18}, and adding it improves the model, \Fstat{1}{180}{12.44}, \plt{.001}, $R^2=.09$ (Table~\ref{tab:a9}, Figure~\ref{fig:a9}). In caf\'es without outdoor seating, the Price Level goes up by 2.74~CZK per 10 seats, \CI{1.41}{4.07}, \plt{.001}. In caf\'es with outdoor seating the slope is only 0.07~CZK, \CI{$-$0.62}{0.75}, \pval{.845}. The Breusch--Pagan test found no unequal variance (\pval{.843}), and the most influential caf\'es are Caf\'e Slavia and Starbucks Reserve (Appendix~\ref{app:diag}). In this sample, size goes with price only in caf\'es without outdoor seating.

\begin{table}[ht]
\centering
\caption{A9: regression of the Price Level (CZK) on the seats inside (per 10 seats), outdoor seating and their interaction, with the slope in each group and the $F$-test of the interaction.}
\label{tab:a9}
\fittable{tables/A9_table.tex}
\end{table}

\begin{figure}[ht]
\centering
\inputfigure[0.7]{figures/A9_figure.pdf}
\caption{A9: Price Level (CZK) against seats inside for caf\'es with ($n=123$) and without ($n=61$) outdoor seating, with one fitted line and 95\% confidence band per group.}
\label{fig:a9}
\end{figure}

\FloatBarrier
\subsection{A10: Does the chance of an above-median Price Level change with size?}
The outcome is 1 if a caf\'e's Price Level is above the median of 79.67~CZK and 0 if not. \Ho: the probability does not change with the seats inside. \Hone: it does change (two-tailed). For 186 caf\'es (3 lost, 91 above the median), the logistic regression gives an average marginal effect of 2.61 percentage points per 10 seats, \CI{$-$0.38}{5.60}, $z=1.71$, \pval{.087}. As a check, the odds ratio per 10 seats is 1.11, \CI{0.98}{1.26} (Table~\ref{tab:a10}, Figure~\ref{fig:a10}). There is no evidence that bigger caf\'es are more likely to be in the more expensive half.

\begin{table}[ht]
\centering
\caption{A10: logistic regression of an above-median Price Level on the seats inside. AME = average marginal effect of 10 more seats, in percentage points (pp); the odds ratio, also per 10 seats, is a check. Events = caf\'es above the median.}
\label{tab:a10}
\fittable{tables/A10_table.tex}
\end{table}

\begin{figure}[ht]
\centering
\inputfigure[0.7]{figures/A10_figure.pdf}
\caption{A10: predicted probability of an above-median Price Level by seats inside, with its 95\% confidence band; dots are the 186 caf\'es (1 = above the median, 0 = not).}
\label{fig:a10}
\end{figure}

\FloatBarrier
\section{Discussion}
Put together, the results show that Prague caf\'es are mostly small places: about 28 seats inside on average, 22 for the median caf\'e, and far fewer outside. Outdoor seating depends on location, and Prague~5 stands out with a median of 0 outdoor seats among the caf\'es we visited. Size matters most for ratings, since small caf\'es are top-rated much more often. It matters little for prices: the link is weak (A6, A7), it is no longer significant once the distance from the Orloj is included (A8), it only shows up in caf\'es without outdoor seating (A9), and it does not change the chance of being in the more expensive half (A10). All ten analyses were planned before looking at the data, so none of them is exploratory.

These conclusions are about this sample, and they have three limits. First, it is a convenience sample: students picked caf\'es they could reach, so some districts (for example Prague~1 and Prague~8) are covered much better than others, and the caf\'es are not a random sample of Prague caf\'es. Second, the data are observational, so none of these links shows a cause. For example, small caf\'es might be rated higher because of the kind of caf\'e they are, not because of their size. Third, caf\'es in the same district are not independent of each other, because they share rents, tourists and competition. All the tests assume independence, so their $p$-values may be a bit too small and their intervals too narrow. The regressions are also affected by a few very big caf\'es, but A7 shows that the direction of the result does not depend on them.

\FloatBarrier
\section{Conclusion}
So how big are Prague caf\'es, and does their size matter? In this sample a caf\'e has about 28 seats inside and 17 more inside than outside, and outdoor seating is rarer in Prague~5. Small caf\'es are rated higher: they are top-rated 28.6 percentage points more often than other caf\'es. Bigger caf\'es are only slightly more expensive, and even that weak link disappears once location is taken into account. Size says more about how a caf\'e is rated than about what it charges.

\appendix

\FloatBarrier
\section{Descriptive statistics}\label{app:desc}
Table~\ref{tab:counts} gives the counts and percentages of the categorical variables, and Table~\ref{tab:desc} the statistics of every numeric variable and every compared group, as produced by \texttt{02\_descriptive.R}.

\begin{table}[ht]
\centering
\caption{Categorical variables: counts and percentages of all 189 caf\'es.}
\label{tab:counts}
\fittable{tables/D_counts.tex}
\end{table}

\begin{table}[ht]
\centering
\caption{Numeric variables, overall and within the compared groups: number of caf\'es, missing values, mean, standard deviation, median, interquartile range and range. Seats in number of seats, Price Level in CZK, distance in km.}
\label{tab:desc}
\fittable{tables/D_stats.tex}
\end{table}

\FloatBarrier
\section{Distributions and normality}\label{app:dist}
Figures~\ref{fig:dhist} to~\ref{fig:dboot} show the histograms, boxplots and Q--Q plots of every variable and group, and the distribution of the bootstrapped means. Table~\ref{tab:dtable} gives the resulting decision. All seat counts are strongly right-skewed (skewness 0.71 to 4.52) because of a few very big caf\'es, and the Price Level and the distance are moderately skewed. None is normal on its own. The bootstrapped means are close to symmetric for every group with at least 30 caf\'es, so the central limit theorem applies there and the parametric method is used for A1, A2, A3 and A6. Prague~2, 3 and 5 have fewer than 30 caf\'es, so A4 uses the non-parametric Kruskal--Wallis test.

\begin{table}[ht]
\centering
\caption{Decision table: for every variable and compared group, the number of caf\'es, the skewness, the skewness of 5{,}000 bootstrapped means (Boot.\ skewness, seed 123), whether the variable is normal ($|$skewness$|<0.5$), whether the central limit theorem applies (CLT: at least 30 caf\'es and $|$boot.\ skewness$|<0.5$), and the resulting method.}
\label{tab:dtable}
\fittable{tables/D_table.tex}
\end{table}

\begin{figure}[ht]
\centering
\inputfigure[1]{figures/D_histograms.pdf}
\caption{Histograms of every variable and compared group (number of caf\'es on the vertical axis; seats, CZK or km on the horizontal axis).}
\label{fig:dhist}
\end{figure}

\begin{figure}[ht]
\centering
\inputfigure[1]{figures/D_boxplots.pdf}
\caption{Boxplots of every variable and compared group: the median (line), the middle 50\% (box) and outliers (dots).}
\label{fig:dbox}
\end{figure}

\begin{figure}[ht]
\centering
\inputfigure[1]{figures/D_qq.pdf}
\caption{Normal Q--Q plots of every variable and compared group. Points on the line mean normal data; points curving away at the ends mean skew.}
\label{fig:dqq}
\end{figure}

\begin{figure}[ht]
\centering
\inputfigure[1]{figures/D_bootstrap.pdf}
\caption{Distribution of 5{,}000 bootstrapped means (seed 123) of every variable and compared group. A symmetric bell shape means the central limit theorem applies.}
\label{fig:dboot}
\end{figure}

\FloatBarrier
\section{Regression diagnostics}\label{app:diag}
Figures~\ref{fig:diag7} to~\ref{fig:diag9} show, for A7, A8 and A9, the residuals against the fitted values (linearity and equal variance), the normal Q--Q plot of the residuals, and the residuals against the leverage with Cook's distance (influential caf\'es). The residuals form a band without a clear funnel, which agrees with the Breusch--Pagan tests ($p=.186$, $.152$ and $.843$), and they are somewhat right-skewed with a few very expensive caf\'es. The most influential caf\'es are VNITROBLOCK (Cook's distance 0.50 in A7), Kav\'arna Creperi\'e (8.1~km from the Orloj, 0.28 in A8) and Caf\'e Slavia (0.31 in A9); none reaches 1. In A7, 11 caf\'es have a Cook's distance above $4/n$; refitting without them changes the slope from 0.66 to 1.92~CZK per 10 seats, so the direction of the result does not depend on them. In A8 the variance inflation factor is 1.01, so seats and distance do not overlap.

\begin{figure}[ht]
\centering
\inputfigure[1]{figures/A7_diagnostics.pdf}
\caption{A7 diagnostics: residuals vs fitted values, normal Q--Q plot of the residuals, and residuals vs leverage with Cook's distance ($n=186$).}
\label{fig:diag7}
\end{figure}

\begin{figure}[ht]
\centering
\inputfigure[1]{figures/A8_diagnostics.pdf}
\caption{A8 diagnostics: residuals vs fitted values, normal Q--Q plot of the residuals, and residuals vs leverage with Cook's distance ($n=186$).}
\label{fig:diag8}
\end{figure}

\begin{figure}[ht]
\centering
\inputfigure[1]{figures/A9_diagnostics.pdf}
\caption{A9 diagnostics: residuals vs fitted values, normal Q--Q plot of the residuals, and residuals vs leverage with Cook's distance ($n=184$).}
\label{fig:diag9}
\end{figure}

\FloatBarrier
\section{Other assumption checks and post hoc tests}\label{app:checks}
\begin{itemize}
  \item A2: Levene's test, $F(1,184)=0.17$, \pval{.680}: the variances of the two groups do not differ. Welch's test was used, which is valid either way.
  \item A4: the full post hoc comparisons are in Table~\ref{tab:a4posthoc}.
  \item A5: the smallest expected count is 19.29, so the chi-square test is valid and Fisher's exact test was not needed.
  \item A7--A9: the Breusch--Pagan tests ($p=.186$, $.152$, $.843$) found no unequal variance, so classic standard errors are reported.
  \item Independence: caf\'es in the same district are not independent, which every test assumes (see the Discussion).
\end{itemize}

\begin{table}[ht]
\centering
\caption{A4 post hoc: Dunn's test for every pair of districts with Holm-adjusted $p$-values, and the difference in medians (second district minus first, seats) with a bootstrapped 95\% confidence interval (5{,}000 resamples, seed 123; not adjusted).}
\label{tab:a4posthoc}
\fittable{tables/A4_posthoc.tex}
\end{table}

\FloatBarrier
\section{Reproducibility}\label{app:repro}
{\sloppy All materials are in the repository \url{https://github.com/sawyer-source/prague-cafe-analysis}. Opening \texttt{prague-cafe-analysis.Rproj} and running the three scripts in order reproduces every number, table and figure of this report from the raw data: \texttt{analysis/01\_cleaning.R} (cleaning, writes \texttt{dataset.csv}), \texttt{analysis/02\_descriptive.R} (description and choice of method) and \texttt{analysis/03\_analysis.R} (analyses A1--A10). The analysis used R 4.3.3 with dplyr 1.1.4, ggplot2 3.4.4, xtable 1.8-4, rstatix 0.7.2, lmtest 0.9-40, sandwich 3.1-0, emmeans 1.10.0 and marginaleffects 0.18.0. Every bootstrap uses 5{,}000 resamples and the seed 123.\par}

\end{document}
